% ============================================================================
% Chapter 10: Using Windows Applications
% ============================================================================

\chapter{Using Windows Applications}
\label{ch:apps}

\begin{objectives}
  \item Open and use built-in Windows applications
  \item Identify the purpose of Notepad, WordPad, Paint, and Calculator
  \item Navigate a basic application window (title bar, menus, toolbar)
  \item Save and close work in an application
\end{objectives}

\bytetip[fun]{Windows comes with free apps already installed --- like getting free games with a new phone! Let's explore the cool tools you already have.}

% ---------------------------------------------------------------
\section{Built-In Windows Applications}
\label{sec:builtin-apps}
\index{Windows!applications}

Windows comes with several useful applications pre-installed. Here are four you'll use often:

% --- Figure 10.1: Applications Overview Cards ---
\begin{figure}[H]
\centering
\begin{tabularx}{\textwidth}{@{}*{4}{>{\centering\arraybackslash}X}@{}}
  \begin{tcolorbox}[colback=ModuleBlue!8, colframe=ModuleBlue, width=3.2cm, height=4.2cm, valign=center, halign=center, arc=8pt]
    {\Huge\textcolor{ModuleBlue}{\faStickyNote}}\vspace{6pt}\\
    \textbf{Notepad}\vspace{3pt}\\
    {\footnotesize Simple text editor\\for quick notes.}\vspace{4pt}\\
    {\scriptsize\itshape Best for: Plain text}
  \end{tcolorbox}
  &
  \begin{tcolorbox}[colback=LabGreen!8, colframe=LabGreen, width=3.2cm, height=4.2cm, valign=center, halign=center, arc=8pt]
    {\Huge\textcolor{LabGreen}{\faFile}}\vspace{6pt}\\
    \textbf{WordPad}\vspace{3pt}\\
    {\footnotesize Basic word processor\\with formatting.}\vspace{4pt}\\
    {\scriptsize\itshape Best for: Short docs}
  \end{tcolorbox}
  &
  \begin{tcolorbox}[colback=FunPurple!8, colframe=FunPurple, width=3.2cm, height=4.2cm, valign=center, halign=center, arc=8pt]
    {\Huge\textcolor{FunPurple}{\faPaintBrush}}\vspace{6pt}\\
    \textbf{Paint}\vspace{3pt}\\
    {\footnotesize Drawing \& basic\\image editing.}\vspace{4pt}\\
    {\scriptsize\itshape Best for: Drawings}
  \end{tcolorbox}
  &
  \begin{tcolorbox}[colback=NoteOrange!8, colframe=NoteOrange, width=3.2cm, height=4.2cm, valign=center, halign=center, arc=8pt]
    {\Huge\textcolor{NoteOrange}{\faCalculator}}\vspace{6pt}\\
    \textbf{Calculator}\vspace{3pt}\\
    {\footnotesize Maths operations\\and conversions.}\vspace{4pt}\\
    {\scriptsize\itshape Best for: Maths}
  \end{tcolorbox}
\end{tabularx}
\caption{Four built-in Windows applications you should know}
\label{fig:app-cards}
\end{figure}

% ---------------------------------------------------------------
\section{Understanding Application Windows}
\label{sec:app-windows}
\index{window!elements}

Every application opens in a \textbf{window}. All windows share common elements:

\begin{theorybox}[Parts of an Application Window]
\begin{itemize}[leftmargin=*, label={\textcolor{ModuleBlue}{\faAngleRight}}, itemsep=3pt]
  \item \textbf{Title Bar} --- Shows the file name and app name; has Minimise, Maximise, and Close buttons
  \item \textbf{Menu Bar} --- Contains menus like File, Edit, Format, View, Help
  \item \textbf{Toolbar / Ribbon} --- Quick-access buttons for common tasks
  \item \textbf{Work Area} --- The main area where you type, draw, or view content
  \item \textbf{Status Bar} --- Shows helpful info at the bottom (zoom, line count, etc.)
\end{itemize}
\end{theorybox}

% --- Figure 10.2: Paint Toolbar Diagram ---
\begin{figure}[H]
\centering
\begin{tikzpicture}
  % Toolbar strip
  \fill[ShadeGray] (0,0) rectangle (14,1.2);
  \draw[NavyBlue, line width=1pt] (0,0) rectangle (14,1.2);
  
  % Sections
  \fill[FunSky!20] (0.2,0.15) rectangle (2.5,1.05);
  \fill[LabGreen!15] (2.7,0.15) rectangle (5.0,1.05);
  \fill[FunPurple!15] (5.2,0.15) rectangle (8.0,1.05);
  \fill[FunYellow!20] (8.2,0.15) rectangle (11.0,1.05);
  \fill[FunCoral!15] (11.2,0.15) rectangle (13.8,1.05);
  
  % Section labels (inside)
  \node[font=\tiny\sffamily\bfseries, text=FunSky!80!black] at (1.35,0.6) {Selection};
  \node[font=\tiny\sffamily\bfseries, text=LabGreen!80!black] at (3.85,0.6) {Drawing};
  \node[font=\tiny\sffamily\bfseries, text=FunPurple!80!black] at (6.6,0.6) {Shapes};
  \node[font=\tiny\sffamily\bfseries, text=FunYellow!80!black] at (9.6,0.6) {Colours};
  \node[font=\tiny\sffamily\bfseries, text=FunCoral!80!black] at (12.5,0.6) {Eraser};
  
  % Labels above with arrows
  \foreach \x/\lbl in {1.35/Selection\\Tools, 3.85/Drawing\\Tools, 6.6/Shape\\Library, 9.6/Colour\\Palette, 12.5/Eraser\\Tool}{
    \draw[NavyBlue, line width=0.8pt, ->] (\x, 1.8) -- (\x, 1.25);
    \node[above, font=\scriptsize\sffamily, text=NavyBlue, align=center] at (\x, 1.8) {\lbl};
  }
\end{tikzpicture}
\caption{The Paint toolbar --- key sections labelled}
\label{fig:paint-toolbar}
\end{figure}

\bytetip[tip]{To open any app quickly, press the \key{Windows} key and start typing the app's name. Windows Search will find it instantly --- no need to click through menus!}

\begin{tryit}
Open \textbf{MS Paint} and draw a simple house with a door, two windows, and a sun in the sky. Use at least \textbf{three different colours}. Save your masterpiece as ``MyHouse.png'' on the Desktop!
\end{tryit}

\begin{funfact}
MS Paint has been included with every version of Windows since \textbf{Windows 1.0 in 1985}. That makes it over 40 years old! Microsoft once threatened to remove it, but there was such a public outcry that they kept it. People \emph{love} Paint!
\end{funfact}

\begin{reviewqs}
  \item Name four built-in Windows applications and describe what each does.
  \item What are the main parts of an application window?
  \item How do you open an application using Windows Search?
  \item What is the difference between Notepad and WordPad?
  \item How do you close an application window?
\end{reviewqs}

\begin{challenge}
Open the \textbf{Calculator} app and use it to solve these problems: (a) 256 $\times$ 48, (b) 7894 $\div$ 23, (c) 15\% of 680. Write down your answers. Then switch to \textbf{Scientific mode} (click the menu $\rightarrow$ Scientific) and calculate $2^{10}$. What do you get?
\end{challenge}
